\section{Forma de Weil, doble círculo y representación espectral}
\label{lectura:manuscrito-03}
La factorización y los relojes de la sección~\ref{lectura:manuscrito-02} producen dos lecturas de una misma aritmética. La primera descompone una forma normal en irreducibles; la segunda registra los retornos de sus acciones multiplicativas. El desarrollo espectral debe conservar ambas procedencias. Su objeto no se obtiene asignando a un retorno un cero conocido de una función: se construyen los funcionales, se establece su compatibilidad y se identifica después su realización analítica.

Este capítulo reúne los desarrollos conservados de la forma completada y del doble círculo. Expone separadamente los resultados incondicionales de normalización, las consecuencias exactas de una identidad positiva de frontera y las realizaciones de esa identidad. Esta separación fija el orden lógico de las demostraciones. En particular, el cálculo de una representación que utiliza la positividad no puede reemplazar la construcción que debe establecerla.

\subsection{Del refinamiento cíclico a la forma completada}
Sea \(N=3^m\), con \(m\geq1\). Sobre el ciclo de \(N\) posiciones se conserva el operador cuadrático del capítulo anterior. Si los índices de Fourier se eligen en el intervalo simétrico

\[
-\frac{N-1}{2}\leq n\leq\frac{N-1}{2},
\]
sus autovalores cuadráticos son

\[
\lambda_{m,n}=\frac{N^2}{\pi}\sin^2\!\left(\frac{\pi n}{N}\right).
\]
La coordenada \(\pi\) es la salida regional anterior del núcleo HMT. En esta etapa interviene en la normalización de la realización analítica, no en la selección de irreducibles. El modo constante se retira y las dos orientaciones de un modo no nulo se cuentan una sola vez:

\[
\Theta_m(x)=\frac12\sum_{n\neq0}e^{-x\lambda_{m,n}}
=\sum_{n=1}^{(N-1)/2}e^{-x\lambda_{m,n}},\qquad x>0.
\]
\HMTresultado{Proposición}{Límite de la traza de calor}{rz-num-03-1}
Para todo \(x>0\),

\[
\lim_{m\to\infty}\Theta_m(x)
=\Theta_+(x):=\sum_{n=1}^{\infty}e^{-\pi n^2x}.
\]
La convergencia es uniforme cuando \(x\) recorre un compacto de \((0,\infty)\).

\textbf{Demostración.} Para un índice fijo \(n\), la equivalencia \(\sin y/y\to1\) da \(\lambda_{m,n}\to\pi n^2\). Además, \(\sin y\geq2y/\pi\) para \(0\leq y\leq\pi/2\), de donde

\[
\lambda_{m,n}\geq\frac{4n^2}{\pi}.
\]
Si \(x\geq a>0\), cada término está dominado por \(e^{-4an^2/\pi}\), cuya suma converge. Extendiendo por cero los términos fuera del intervalo de índices del ciclo y aplicando convergencia dominada se obtienen el límite y su uniformidad en compactos. \hfill\(\square\)

La inversión de esta traza explica por qué la completación contiene una contribución polar además de los modos positivos. Defínase

\[
\vartheta(x)=1+2\Theta_+(x).
\]
La periodización de la gaussiana

\[
F_x(t)=\sum_{n\in\mathbb Z}e^{-\pi x(t+n)^2}
\]
converge con todas sus derivadas. Su coeficiente de Fourier de índice \(k\) es \(x^{-1/2}e^{-\pi k^2/x}\), obtenido integrando la gaussiana sobre la recta. La serie de Fourier converge absolutamente; al evaluarla en \(t=0\) se obtiene

\[
\vartheta(x)=x^{-1/2}\vartheta(1/x).
\]
En términos de la traza positiva,

\[
\Theta_+(x)=x^{-1/2}\Theta_+(1/x)+\frac{x^{-1/2}-1}{2}.
\]
El último sumando conserva la diferencia entre retirar el modo constante antes o después de invertir la escala.

\HMTresultado{Proposición}{Transformada de Mellin y continuación}{rz-num-03-2}
Para \(\Re s>1\), la función de partición aritmética \(Z(s)=\sum_{n\geq1}n^{-s}\) satisface

\[
\Lambda(s):=\int_0^\infty\Theta_+(x)x^{s/2-1}\,dx
=\pi^{-s/2}\Gamma(s/2)Z(s).
\]
La función

\[
\xi(s)=\frac12s(s-1)\Lambda(s)
\]
admite una prolongación entera y satisface \(\xi(s)=\xi(1-s)\).

\textbf{Demostración.} La convergencia absoluta para \(\Re s>1\) permite integrar término a término. El cambio \(u=\pi n^2x\) da el factor \(\pi^{-s/2}\Gamma(s/2)n^{-s}\). Para continuar la integral, se separan los intervalos \((0,1)\) y \((1,\infty)\) y se aplica la inversión anterior al primero. Resulta

\[
\Lambda(s)=\frac1{s(s-1)}
+\int_1^\infty\Theta_+(x)
\left(x^{s/2-1}+x^{(1-s)/2-1}\right)\,dx.
\]
La integral es entera en \(s\): la disminución exponencial de \(\Theta_+(x)\) domina uniformemente cualquier potencia y cualquier potencia de \(\log x\) sobre compactos del plano. El factor \(s(s-1)\) elimina los dos polos mostrados explícitamente. La expresión es invariante bajo \(s\mapsto1-s\), lo que prueba la ecuación funcional. \hfill\(\square\)

La transformada de Mellin de la traza no es \(Z(s)\) sin más: contiene los factores \(\pi^{-s/2}\Gamma(s/2)\). Mantenerlos explícitos es indispensable, pues sus contribuciones reaparecen en la forma cuadrática completada. El producto por irreducibles del capítulo anterior y esta construcción aditiva coinciden en el semiplano de convergencia por una identidad demostrada de funcionales.

\subsection{Las contribuciones aritmética, gamma y polar}
\label{rz:forma}
Fíjese \(R>0\) y considérese

\[
\mathcal D_R=C_c^\infty((-R/2,R/2)).
\]
El producto interior será lineal en la primera variable. Se utilizan

\[
\widehat f(z)=\int_{\mathbb R}f(x)e^{-izx}\,dx,\qquad
(T_af)(x)=f(x-a),
\]
\[
C_{f,g}(a)=\langle f,T_ag\rangle_{L^2(\mathbb R)},\qquad
h_{f,g}(z)=\widehat f(z)\overline{\widehat g(\overline z)}.
\]
Estas definiciones fijan simultáneamente la orientación de las traslaciones y la normalización de Fourier. Para \(f,g\in\mathcal D_R\), \(C_{f,g}(a)=0\) si \(|a|\geq R\).

Los irreducibles ya construidos suministran los índices primo–potencia y sus pesos:

\[
\ell_{p,k}=k\log p,\qquad
w_{p,k}=(\log p)p^{-k/2},\qquad k\geq1.
\]
La contribución procedente de la completación gamma utiliza

\[
\mu(a)=\frac{e^{-a/2}}{1-e^{-2a}},\qquad
c_\Gamma=\psi(1/4)-\log\pi,
\]
donde \(\psi=\Gamma'/\Gamma\). En particular, \(c_\Gamma<0\). Por ejemplo, la monotonía de \(\psi\) en el eje positivo, obtenida de \(\psi'(x)=\sum_{n\geq0}(n+x)^{-2}>0\), y \(\psi(1)=-\gamma_{\mathrm E}<0\) dan esta desigualdad.

La forma completada es

\[
\begin{aligned}
\mathscr W_R(f,g)={}&h_{f,g}(i/2)+h_{f,g}(-i/2)
+c_\Gamma C_{f,g}(0)\\
&+\int_0^\infty\mu(a)
\bigl(2C_{f,g}(0)-C_{f,g}(a)-C_{f,g}(-a)\bigr)\,da\\
&-\sum_{\ell_{p,k}\leq R}w_{p,k}
\bigl(C_{f,g}(\ell_{p,k})+C_{f,g}(-\ell_{p,k})\bigr).
\end{aligned}
\]
La suma primo–potencia es finita: sólo intervienen \(p^k\leq e^R\). No se ha sustituido la forma por un muestreo finito de ceros.

\HMTresultado{Proposición}{Existencia y compatibilidad de la forma}{rz-num-03-3}
La expresión anterior define una forma sesquilineal hermítica y continua en \(\mathcal D_R\). Si \(R<R'\), sus restricciones coinciden:

\[
\mathscr W_{R'}|_{\mathcal D_R\times\mathcal D_R}=\mathscr W_R.
\]
\textbf{Demostración.} Como las traslaciones son unitarias,

\[
2C_{f,g}(0)-C_{f,g}(a)-C_{f,g}(-a)
=\langle(I-T_a)f,(I-T_a)g\rangle.
\]
Para \(a\to0\), cada diferencia tiene norma \(O(a)\), mientras que \(\mu(a)=O(a^{-1})\). El integrando es, por tanto, \(O(a)\). Para \(a\to\infty\), las diferencias tienen normas acotadas y \(\mu(a)=O(e^{-a/2})\). La integral converge absolutamente. Estas mismas estimaciones, expresadas mediante las normas de \(f,g\) y sus primeras derivadas sobre un compacto fijo, prueban la continuidad. La simetría hermítica resulta de intercambiar \(f,g\) y conjugar. Si se amplía \(R\), los términos primo–potencia añadidos tienen correlaciones nulas sobre \(\mathcal D_R\), lo que prueba la compatibilidad. \hfill\(\square\)

Puede escribirse así una única forma \(\mathscr W\) sobre \(\mathcal D=C_c^\infty(\mathbb R)=\bigcup_R\mathcal D_R\). La compatibilidad corresponde a la diferencia completa de contribuciones; no exige que cada columna auxiliar conserve por separado su norma al ampliar el soporte.

\subsection{Columnas de transporte e identidad positiva de frontera}
\label{rz:columnas}
La representación de defecto uniforme del capítulo anterior da, para un ciclo de longitud \(N\),

\[
\delta_N(T_a)^*\delta_N(T_b)=(I-T_a)^*(I-T_b).
\]
Esta igualdad permite realizar las diferencias de traslación en una hoja de transporte normalizada. Para exponer el argumento sin repetir esa realización, se escribirá simplemente \(\delta_af=(I-T_a)f\). Las evaluaciones polares se descomponen en

\[
b_\pm(f)=\frac{\widehat f(i/2)\pm\widehat f(-i/2)}{\sqrt2}.
\]
Defínanse dos aplicaciones sobre \(\mathcal D_R\):

\[
J_{+,R}f=
\left(
(\sqrt{w_{p,k}}\,\delta_{\ell_{p,k}}f)_{\ell_{p,k}\leq R},
(\sqrt{\mu(a)}\,\delta_af)_{a>0},
b_+(f)
\right),
\]
\[
J_{-,R}f=
\left(
(\sqrt{2w_{p,k}}\,f)_{\ell_{p,k}\leq R},
\sqrt{-c_\Gamma}\,f,
b_-(f)
\right).
\]
Los espacios de llegada son sumas ortogonales de copias de \(L^2(\mathbb R)\), una integral directa en la componente gamma y una componente escalar. Cuando se usen sus imágenes como espacios efectivos, se tomarán los cierres de \(\operatorname{ran}J_{+,R}\) y \(\operatorname{ran}J_{-,R}\). El cierre de una imagen diagonal no debe confundirse con la suma de todas las coordenadas del espacio ambiente.

\HMTresultado{Proposición}{Descomposición exacta de la forma}{rz-num-03-4}
Se cumple

\[
\mathscr W_R(f,g)
=\langle J_{+,R}f,J_{+,R}g\rangle
-\langle J_{-,R}f,J_{-,R}g\rangle.
\]
\textbf{Demostración.} Para cada primo–potencia, la primera columna aporta

\[
w_{p,k}\bigl(2C_{f,g}(0)
-C_{f,g}(\ell_{p,k})-C_{f,g}(-\ell_{p,k})\bigr);
\]
la segunda retira \(2w_{p,k}C_{f,g}(0)\). La componente gamma es precisamente la integral de la proposición~\ref{rz-num-03-3}. La componente constante negativa produce \(c_\Gamma C_{f,g}(0)\). Por último,

\[
b_+(f)\overline{b_+(g)}-b_-(f)\overline{b_-(g)}
=h_{f,g}(i/2)+h_{f,g}(-i/2).
\]
La suma de estas identidades da la expresión anunciada. \hfill\(\square\)

Esta descomposición no establece por sí sola positividad: es una diferencia de dos formas positivas. La identidad de frontera del propietario demostrativo utiliza un espacio \(\mathcal K_R\), una proyección ortogonal \(P_R\) y un codificador \(\Xi_R:\mathcal D_R\to\mathcal K_R\) cuyos momentos deben satisfacer

\[
\langle\Xi_Rf,\Xi_Rg\rangle
=\langle J_{+,R}f,J_{+,R}g\rangle,
\]
\[
\langle P_R\Xi_Rf,P_R\Xi_Rg\rangle
=\langle J_{-,R}f,J_{-,R}g\rangle.
\]
\HMTresultado{Teorema}{Consecuencia de los dos momentos de frontera}{rz-num-03-5}
Si las dos identidades anteriores se cumplen sobre \(\mathcal D_R\), la aplicación \(E_R=(I-P_R)\Xi_R\) satisface

\[
\mathscr W_R(f,g)=\langle E_Rf,E_Rg\rangle.
\]
En particular, \(\mathscr W_R(f,f)\geq0\).

\textbf{Demostración.} La descomposición ortogonal de \(\Xi_Rf\) mediante \(P_R\) y \(I-P_R\) da

\[
\langle\Xi_Rf,\Xi_Rg\rangle
=\langle P_R\Xi_Rf,P_R\Xi_Rg\rangle
+\langle E_Rf,E_Rg\rangle.
\]
Se sustituyen los dos momentos y se aplica la proposición~\ref{rz-num-03-4}. \hfill\(\square\)

La importancia del teorema reside en la identificación de la forma aritmética con una norma de frontera. El paso que determina su aplicación es la acción del codificador y la prueba de ambos momentos; la ortogonalidad de una proyección no los impone para una forma arbitraria. El apéndice de procedencia de esta edición registra el propietario conservado y el punto de su integración material. No se sustituye esa acción por un operador definido retrospectivamente a partir de la positividad que se quiere establecer.

Hay una segunda conservación, útil una vez construido \(E_R\). La matriz de incidencia

\[
B=\begin{pmatrix}7&2\\2&7\end{pmatrix}=9P_++5P_-
\]
tiene proyectores ortogonales

\[
P_\pm=\frac12
\begin{pmatrix}1&\pm1\\\pm1&1\end{pmatrix}.
\]
Su normalización da \(A=P_++qP_-\), con \(q=5/9\). Si \(D=\sqrt{1-q^2}\,P_-\), entonces \(A^*A+D^*D=I\). Por iteración,

\[
\|x\|^2=\sum_{n=0}^{M-1}\|DA^nx\|^2+\|A^Mx\|^2.
\]
Sobre la hoja \(P_-x=x\), el término terminal es \(q^{2M}\|x\|^2\); por consiguiente, la aplicación

\[
Lx=(DA^nx)_{n\geq0}
\]
es isométrica en esa hoja. La telescopía explica la conservación mediante la distribución de pérdidas; la pertenencia de \(E_Rf\) a la hoja debe proceder de la construcción anterior. Extinguir un término terminal no selecciona retrospectivamente la forma de la que partía.

\subsubsection*{Codificador interno de las historias y sus dos momentos finitos}
El transporte correlativo del TPK ya permite construir un codificador antes de definir su residuo. Fíjese un horizonte \(N\geq1\). Para cada nivel, sea \(\mathscr H_k\) la suma ortogonal de las dos fibras de hoja sobre las historias completas de ese nivel. Si \(d(x)\) cuenta las emisiones salientes de una historia \(x\), el transporte actúa por

\[
U_k\iota_xv=
\sum_{\alpha:x\to y}\frac1{\sqrt{d(x)}}\,\iota_yU_\alpha v.
\]
Se considera aquí la realización isométrica del transporte de hoja: \(U_\alpha^*U_\alpha=I\). Historias distintas tienen descendientes distintos porque se conserva la ruta completa. Por tanto,

\[
\left.U_k^*U_k\right|_{\mathscr H_x}
=\frac1{d(x)}\sum_{\alpha:x\to y}U_\alpha^*U_\alpha=I.
\]
La consistencia de los pesos y la conservación de la memoria se desarrollan en la sección~\ref{lectura:manuscrito-05} y se aplican al transporte espectral en la subsección~\ref{rz:refinamiento}. La isometría de cada transporte de hoja es una condición suficiente explícita para esta igualdad; la normalización de los pesos, aisladamente, no la demuestra.

Sean \(P_k,Q_k\) los proyectores ortogonales de las hojas, con \(P_k+Q_k=I\). Las partes del transporte que conservan y cambian hoja son

\[
A_k=P_{k+1}U_kP_k+Q_{k+1}U_kQ_k,
\qquad
\eta_k=Q_{k+1}U_kP_k+P_{k+1}U_kQ_k.
\]
El cálculo de bloques da exactamente

\[
A_k^*A_k+\eta_k^*\eta_k
=P_kU_k^*U_kP_k+Q_kU_k^*U_kQ_k=I.
\]
La primera igualdad es válida para cualquier transporte acotado. La segunda utiliza la realización isométrica indicada. En general, el primer miembro es la parte diagonal por hojas del Gram de \(U_k\), no necesariamente su Gram completo.

Defínanse \(F_{0,0}=I\), \(F_{0,k}=A_{k-1}\cdots A_0\) para \(k\geq1\), y \(T_N=F_{0,N}^*F_{0,N}\). El codificador interno tiene dominio y acción

\[
\Xi_N^{\mathrm{int}}:\mathscr H_0\longrightarrow
\mathscr K_N^{\mathrm{int}}
:=\left(\bigoplus_{k=0}^{N-1}\mathscr H_{k+1}\right)\oplus\mathscr H_N,
\]
\[
\Xi_N^{\mathrm{int}}v
=\bigl(\eta_kF_{0,k}v\bigr)_{0\leq k<N}\oplus F_{0,N}v.
\]
El último sumando se mantiene separado incluso cuando coincide como espacio con el de una salida anterior. Sea \(\Pi_N^{\mathrm{term}}\) la proyección ortogonal sobre ese sumando terminal. Se obtienen los dos momentos internos

\[
(\Xi_N^{\mathrm{int}})^*\Xi_N^{\mathrm{int}}=I,
\qquad
(\Xi_N^{\mathrm{int}})^*\Pi_N^{\mathrm{term}}\Xi_N^{\mathrm{int}}=T_N.
\]
\textbf{Demostración.} La identidad de bloques implica

\[
F_{0,k}^*F_{0,k}-F_{0,k+1}^*F_{0,k+1}
=F_{0,k}^*\eta_k^*\eta_kF_{0,k}.
\]
Sumar desde \(k=0\) hasta \(N-1\) produce \(I=\sum_{k<N}F_{0,k}^*\eta_k^*\eta_kF_{0,k}+T_N\), que es el primer Gram anunciado. La proyección terminal retiene solamente \(F_{0,N}v\), por lo que su Gram es \(T_N\). Así, el complemento \((I-\Pi_N^{\mathrm{term}})\Xi_N^{\mathrm{int}}\) registra las salidas de hoja y tiene Gram \(I-T_N\). Todos estos operadores proceden de la misma colección de transportes; el codificador no se define a partir de un residuo cuya positividad se suponga previamente. \hfill\(\square\)

Este resultado recupera un codificador efectivo con sus dos momentos, dentro del constructor conjunto que conserva las otras cuatro lecturas del continuo. Su aplicación a las columnas completadas de esta sección tiene un contenido adicional preciso: habría que identificar una acción previa \(R_R:\mathcal D_R\to\mathscr H_0\), con la compatibilidad de horizontes que exija el dominio de prueba, para la cual

\[
\langle R_Rf,R_Rg\rangle
=\langle J_{+,R}f,J_{+,R}g\rangle,
\]
\[
\langle F_{0,N}R_Rf,F_{0,N}R_Rg\rangle
=\langle J_{-,R}f,J_{-,R}g\rangle.
\]
La composición \(\Xi_N^{\mathrm{int}}R_R\) tendría entonces los momentos del teorema~\ref{rz-num-03-5}. Esta identificación de \(R_R\) con las columnas primo–potencia, gamma y polar no está reunida aquí. La construcción interna anterior sí lo está, con su acción y prueba; no se la presenta como demostración de la positividad global de \(\mathscr W\). Tampoco se suprime \(T_N\), ni se afirma que su límite sea cero.

\subsubsection*{Compatibilidad del horizonte y tipo de la proyección}
La proyección terminal \(\Pi_N^{\mathrm{term}}\) selecciona el último sumando del codificador de historias. La expectativa uniforme \(P_R\) del teorema~\ref{rz-num-03-5} selecciona, en cambio, el componente uniforme de un calendario de nueve posiciones. La igualdad de sus funciones en una composición debe demostrarse mediante la aplicación concreta entre esos espacios; ambas son proyecciones ortogonales, pero esa propiedad común no identifica sus imágenes.

La dependencia respecto del horizonte impone asimismo una identidad comprobable. Sean \(R\) una misma aplicación y \(J_-\) una columna fija. Si, para dos horizontes consecutivos, se exige

\[
R^*T_NR=J_-^*J_-=R^*T_{N+1}R,
\]
entonces

\[
\eta_NF_{0,N}R=0.
\]
\textbf{Demostración.} Al restar las igualdades y utilizar la identidad telescópica se obtiene

\[
0=R^*(T_N-T_{N+1})R
=(\eta_NF_{0,N}R)^*(\eta_NF_{0,N}R).
\]
Su evaluación en cada vector implica la anulación indicada. Si la igualdad de momentos se mantiene con el mismo \(R\) en una sucesión ilimitada de horizontes y \(T_N\) tiende fuertemente a cero, se obtiene además \(J_-=0\). \hfill\(\square\)

Por tanto, una identificación con columna negativa no nula debe especificar su horizonte o su colección compatible de aplicaciones \(R_{R,N}\), o conservar el terminal límite que realmente corresponda. La contracción de un residuo de salida no puede trasladarse sin esa distinción a la columna completa que representa el segundo momento. Esta precisión mantiene íntegro el codificador recuperado y determina qué compatibilidad adicional debe verificar su realización analítica.

\subsection{Representación de traslaciones}
\label{rz:traslaciones}
Supóngase establecida sobre todo \(\mathcal D\) una identidad positiva

\[
\mathscr W(f,g)=\langle Ef,Eg\rangle
\]
compatible con los dominios de soporte. El espacio espectral se construye a partir de la propia forma:

\[
\mathcal N=\{f\in\mathcal D:\mathscr W(f,f)=0\},\qquad
\mathcal H_{\mathscr W}
=\overline{\mathcal D/\mathcal N}.
\]
La desigualdad de Cauchy–Schwarz para una forma positiva prueba que el producto interior del cociente está bien definido. La aplicación \([f]\mapsto Ef\) se extiende a un isomorfismo unitario desde \(\mathcal H_{\mathscr W}\) sobre \(\overline{\operatorname{ran}E}\).

\HMTresultado{Proposición}{Evolución unitaria y generador}{rz-num-03-6}
Las traslaciones inducen un grupo unitario fuertemente continuo

\[
U_t[f]=[T_tf].
\]
Su generador autoadjunto \(H\), con la convención \(U_t=e^{itH}\), actúa en las clases de funciones de prueba por

\[
H[f]=[i f'].
\]
\textbf{Demostración.} Las correlaciones son invariantes bajo la traslación común de sus dos argumentos. Además, \(\widehat{T_tf}(z)=e^{-izt}\widehat f(z)\); en \(h_{T_tf,T_tg}(z)\) los factores se cancelan también para \(z\) complejo. Por ello \(\mathscr W(T_tf,T_tg)=\mathscr W(f,g)\), y la traslación preserva \(\mathcal N\). Define una isometría sobre el cociente, invertida por \(U_{-t}\), que se prolonga al espacio completo.

Para \(t\to0\), \(T_tf-f\) tiende a cero en todas las seminormas de funciones de prueba sobre un compacto común. La continuidad probada para \(\mathscr W\) implica continuidad fuerte sobre el cociente denso; la unitariedad la extiende a todo el espacio. El teorema de Stone proporciona un generador autoadjunto. La derivada en \(t=0\) es \([-f']\), que coincide con \(iH[f]\). \hfill\(\square\)

Las funciones de prueba forman además un núcleo del generador. Puede verse regularizando un vector mediante \(\int\chi(t)U_t\,dt\), con \(\chi\in C_c^\infty(\mathbb R)\). Esta regularización y su derivada son operadores acotados, controlados por \(\|\chi\|_1\) y \(\|\chi'\|_1\). Se aproximan primero los vectores regularizados por clases de funciones de prueba y luego se aproxima la identidad con una colección de regularizadores. La aproximación obtenida es en la norma del grafo \(\|x\|+\|Hx\|\), lo que prueba la afirmación.

La representación se ha construido a partir de una forma positiva y de sus traslaciones. Este argumento no utiliza como entradas una lista de ordenadas de ceros ni una asignación de ceros a vueltas del calendario.

\subsection{Identificación analítica y multiplicidades}
\label{rz:identificacion}
La forma explícita anterior se reconoce como la forma completada de Weil. Con

\[
\Phi_f(s)=\widehat f\bigl(i(s-1/2)\bigr),
\qquad \rho^\sharp=1-\overline\rho,
\]
la fórmula explícita proporciona su expresión espectral

\[
\mathscr W(f,g)=\sum_\rho
\Phi_f(\rho)\overline{\Phi_g(\rho^\sharp)},
\]
con los ceros no triviales de la función completada y sus multiplicidades.
La formulación global del criterio de Weil afirma que la positividad para
todas las funciones de prueba equivale a que \(\Re\rho=1/2\) para todos
esos ceros. Se utiliza el criterio en la formulación de
Bombieri~\cite{bombieri2000}. La condición global no puede reemplazarse
por positividad de un número finito de matrices.

\subsubsection*{Diccionario de pruebas Mellin--Fourier}
\label{rz:mellin-fourier}
Para explicitar la normalización de ese reconocimiento, sea
\(\mathcal D=C_c^\infty(\mathbb R;\mathbb C)\). La aplicación
\[
 (\mathcal T_{\log}f)(x)=x^{-1/2}f(\log x),\qquad x>0,
\]
es una biyección sobre \(C_c^\infty((0,\infty);\mathbb C)\), con inversa
\[
 (\mathcal T_{\log}^{-1}u)(t)=e^{t/2}u(e^t).
\]
En efecto, el cambio exponencial lleva los soportes compactos a compactos
separados de cero y de infinito. Las fórmulas son inversas y sus factores
son suaves. Definiendo
\(\widehat u_{\rm Mel}(s)=\int_0^\infty u(x)x^{s-1}\,dx\),
la sustitución \(x=e^t\) da, para todo \(s\in\mathbb C\),
\[
 \widehat{\mathcal T_{\log}f}_{\rm Mel}(s)
 =\int_{\mathbb R}f(t)e^{(s-1/2)t}\,dt
 =\widehat f\bigl(i(s-1/2)\bigr)=\Phi_f(s).
\]
La conjugación se conserva exactamente:
\[
 \widehat{\overline{\mathcal T_{\log}g}}_{\rm Mel}(1-s)
 =\overline{\Phi_g(1-\overline s)}.
\]
Así, los dos factores del sumando espectral coinciden con los de la
fórmula explícita en coordenadas multiplicativas, sin factor residual
y contando las mismas multiplicidades. La suma converge absolutamente:
la integración por partes da decaimiento rápido de las transformadas
en bandas verticales y el recuento de ceros es \(O(T\log T)\).
La biyección traslada el cuantificador sobre todas las pruebas, no una
selección de ellas. La fórmula explícita y la implicación difícil del
criterio son los antecedentes citados; el cambio de variables anterior
demuestra su correspondencia de dominios y normalizaciones, no una nueva
prueba de positividad. Este dominio compacto no se declara estable bajo
el Cayley integral: su realización posterior usa el espacio exponencial
de la subsección~\ref{tm-add-m4}.

Una vez satisfechos esos antecedentes, escribiendo \(\rho=1/2+i\gamma\), resulta

\[
\mathscr W(f,g)
=\sum_{\gamma\in\Gamma_\xi}
m_\gamma\widehat f(-\gamma)\overline{\widehat g(-\gamma)}.
\]
Aquí \(\Gamma_\xi\) contiene las ordenadas distintas y \(m_\gamma\) es la multiplicidad del cero correspondiente. Sea

\[
\mu_\xi=\sum_{\gamma\in\Gamma_\xi}m_\gamma\delta_\gamma.
\]
La aplicación

\[
\mathcal F_\xi[f](\gamma)=\widehat f(-\gamma)
\]
es isométrica de \(\mathcal H_{\mathscr W}\) en \(L^2(\Gamma_\xi,\mu_\xi)\). En ella, \(U_t\) actúa como multiplicación por \(e^{it\gamma}\) y \(H\) como multiplicación por \(\gamma\).

\HMTresultado{Proposición}{Exhaustividad de la representación atómica}{rz-num-03-7}
La aplicación \(\mathcal F_\xi\) es sobreyectiva.

\textbf{Demostración.} Su imagen cerrada es invariante bajo las multiplicaciones por \(e^{it\gamma}\) para \(t\) positivo y negativo; por tanto es reductora. Para una ordenada fija \(\gamma_0\), las medias

\[
\frac1{2T}\int_{-T}^T e^{-it\gamma_0}U_t\,dt
\]
convergen fuertemente al proyector sobre el átomo \(\gamma_0\). Esto se comprueba primero coordenada a coordenada y después por convergencia dominada en el espacio de sucesiones. Existe \(f\in C_c^\infty\) con \(\widehat f(-\gamma_0)\neq0\); por ejemplo, basta modular una función de integral no nula. La imagen contiene así el vector de ese átomo. Como contiene todos los vectores de soporte finito y es cerrada, coincide con el espacio completo. \hfill\(\square\)

La medida conserva la multiplicidad aritmética como peso. Conviene precisar su relación con la multiplicidad de un operador: el espacio escalar \(L^2(\Gamma_\xi,\mu_\xi)\) tiene un subespacio propio de dimensión uno por cada átomo, aunque \(m_\gamma>1\). Por ello, si una función \(F\) produce un operador de clase traza,

\[
\operatorname{Tr}(F(H))=\sum_{\gamma\in\Gamma_\xi}F(\gamma),
\]
mientras que la suma que conserva los pesos de los ceros es

\[
\int F(\gamma)\,d\mu_\xi(\gamma)
=\sum_{\gamma\in\Gamma_\xi}m_\gamma F(\gamma).
\]
Esta segunda suma puede representarse como una traza ponderada \(\operatorname{Tr}(M_mF(H))\), cuando el producto sea de clase traza, con \(M_m\) el operador de multiplicación por \(m_\gamma\). Otra construcción, distinta, sería una fibra de dimensión \(m_\gamma\) en cada átomo. No se intercambian esas dos representaciones sin declararlo.

\subsection{Los dos círculos y la orientación espectral}
\label{rz:circulos}
La primera circunferencia es una coordenada compacta de la recta espectral. Para un número real \(\gamma\), defínase

\[
\tau_\gamma=\frac{\gamma+i/2}{\gamma-i/2}.
\]
Su módulo es uno y \(\tau_\gamma\neq1\). Las fórmulas inversas son

\[
\gamma=\frac{i}{2}\frac{\tau_\gamma+1}{\tau_\gamma-1},
\qquad
\gamma=\frac12\cot\frac{\theta}{2}
\quad\text{si}\quad\tau_\gamma=e^{i\theta}.
\]
Si \(s=1/2+i\gamma\), también

\[
\tau_\gamma=\frac{s-1}{s},\qquad s=\frac1{1-\tau_\gamma}.
\]
No se ha perdido la ordenada al pasar a la fase, siempre que se mantenga el punto de la circunferencia y la convención de orientación.

Para el generador autoadjunto anterior, el cálculo funcional define el operador unitario

\[
C=(H+iI/2)(H-iI/2)^{-1}.
\]
La igualdad \(C^*C=I\) se deduce de que los dos factores son funciones del mismo operador autoadjunto y \(|(\gamma+i/2)/(\gamma-i/2)|=1\). El punto \(1\) no corresponde a una ordenada finita, pero puede pertenecer al espectro de \(C\) como punto de acumulación cuando \(H\) es no acotado.

La segunda circunferencia es el dual de la memoria entera: a una fase \(\tau\) se asocia el carácter

\[
\chi_\tau(n)=\tau^n,\qquad n\in\mathbb Z.
\]
La conexión nonádica devuelve la fase visible a su posición inicial y aumenta en una unidad la memoria de una vuelta. La realización espectral asigna a esa unidad la acción de \(C\); sobre un vector propio de ordenada \(\gamma\), la memoria de \(n\) vueltas actúa por \(\tau_\gamma^n\). El primer círculo codifica la ordenada; el segundo expresa cómo su fase actúa sobre una memoria aditiva. Son dos funciones matemáticas relacionadas, no dos listas que deban identificarse por su posición ordinal.

\subsection{Refinamiento cilíndrico con memoria}
\label{rz:refinamiento}
La conexión debe actuar sobre la estructura que se refina. En el árbol común del TPK, un vértice es una historia finita de emisiones admisibles; dos historias que llegan al mismo estado terminal siguen siendo distintas cuando sus registros de transporte difieren. Si \(\operatorname{Out}(x)\) es el conjunto finito de emisiones salientes, contadas con su multiplicidad, se escribe

\[
d(x)=|\operatorname{Out}(x)|\geq1.
\]
La prolongación neutra conserva la no vacuidad de los niveles. El peso local declarado por el árbol es \(1/d(x)\). Para una historia \(h=(\epsilon_1,\ldots,\epsilon_k)\), con estados sucesivos \(x_0,\ldots,x_k\), su peso es

\[
\mu_k(h)=\prod_{j=0}^{k-1}\frac1{d(x_j)}.
\]
La suma de los pesos de todos sus sucesores es

\[
\sum_{\epsilon\in\operatorname{Out}(x_k)}
\mu_{k+1}(h\epsilon)
=\mu_k(h)\sum_{\epsilon\in\operatorname{Out}(x_k)}\frac1{d(x_k)}
=\mu_k(h).
\]
Partiendo de la raíz de peso uno, esta identidad demuestra por inducción la normalización y la consistencia proyectiva. Si una simetría reversible transporta biyectivamente las emisiones salientes con sus multiplicidades, conserva \(d(x)\), todos estos pesos y sus medidas de cilindros. Así, la medida utilizada a continuación procede de la regla de ramificación declarada, antes de asignar un operador a sus fibras. No se confunde con una distribución uniforme sobre estados terminales que identificara historias diferentes.

En forma general, sea \(\mathcal Z_k\) un nivel finito de estados cilíndricos y sea \(\pi_k:\mathcal Z_{k+1}\to\mathcal Z_k\) su mapa de truncamiento. Supóngase dada una colección de pesos positivos compatibles,

\[
\mu_k(z)=\sum_{\pi_k(z')=z}\mu_{k+1}(z'),
\qquad \sum_{z\in\mathcal Z_k}\mu_k(z)=1.
\]
Los pesos pertenecen a la medida de estados que se transporta. La forma de un cilindro numérico visible no basta por sí sola para reconstruirlos.

Sean \(q_k:\mathcal Z_k\to\mathbb Z\) las coordenadas de memoria. El incremento de una prolongación es

\[
m(z',z)=q_{k+1}(z')-q_k(z).
\]
Al concatenar prolongaciones se verifica exactamente la ley de cociclo:

\[
m(z'',z)=m(z'',z')+m(z',z).
\]
Defínanse

\[
w(z',z)=
\sqrt{\frac{\mu_{k+1}(z')}{\mu_k(z)}},
\qquad
\mathcal K_k=\ell^2(\mathcal Z_k)\otimes\mathcal H_{\mathscr W}.
\]
El transporte con memoria es

\[
\widetilde U_k(e_z\otimes v)
=\sum_{\pi_k(z')=z}w(z',z)e_{z'}\otimes C^{m(z',z)}v.
\]
\HMTresultado{Proposición}{Isometría y composición del transporte}{rz-num-03-8}
Cada \(\widetilde U_k\) es isométrico. La composición sobre varios niveles tiene la misma forma, con el cociente de pesos entre los niveles extremos y el incremento total de memoria.

\textbf{Demostración.} Dos predecesores distintos tienen conjuntos disjuntos de descendientes. Para un predecesor fijo, las potencias de \(C\) son unitarias y \(\sum_{z'\mapsto z}w(z',z)^2=1\). Éstas son exactamente las identidades que conservan productos interiores de los vectores de base. En una composición, el producto de pesos telescopa:

\[
\sqrt{\frac{\mu_{k+2}(z'')}{\mu_{k+1}(z')}}
\sqrt{\frac{\mu_{k+1}(z')}{\mu_k(z)}}
=\sqrt{\frac{\mu_{k+2}(z'')}{\mu_k(z)}}.
\]
Las potencias de \(C\) suman sus exponentes por la ley de cociclo. La repetición prueba la afirmación para cualquier número finito de niveles. \hfill\(\square\)

La misma identidad admite una lectura por cambio de representación. Si \(U_k\) denota el transporte sin memoria y

\[
G_k(e_z\otimes v)=e_z\otimes C^{q_k(z)}v,
\]
entonces

\[
\widetilde U_k=G_{k+1}(U_k\otimes I)G_k^*.
\]
No se pierde la coordenada de memoria al formular la isometría: aparece en los dos extremos del transporte. Esta igualdad también muestra que la representación torcida utiliza un operador \(C\) ya especificado. Por sí sola no identifica cualquier transporte nativo previo con ese operador; la identificación exige componer sus mapas y comprobar sus acciones sobre el mismo dominio.

\HMTresultado{Proposición}{Retorno nonádico sobre la sección compatible}{rz-num-03-9}
Defínanse

\[
\Omega_k=\sum_z\sqrt{\mu_k(z)}e_z,\qquad
J_kv=\Omega_k\otimes v.
\]
Si toda trayectoria del retorno de \(k\) a \(k+9\) tiene incremento de memoria uno, entonces

\[
\widetilde U_{\Gamma,k}J_k=J_{k+9}C.
\]
Para \(r\) vueltas consecutivas,

\[
\widetilde U_{\Gamma,k+9(r-1)}\cdots
\widetilde U_{\Gamma,k}J_k=J_{k+9r}C^r.
\]
\textbf{Demostración.} Al multiplicar el peso de cada descendiente por el coeficiente \(\sqrt{\mu_k(z)}\) de su predecesor queda \(\sqrt{\mu_{k+9}(z')}\). La condición sobre el incremento permite sacar el mismo operador \(C\) de cada sumando. Se obtiene la primera igualdad. La segunda se sigue por inducción. \hfill\(\square\)

Los índices de nivel son esenciales. El retorno es de la fase y la sección se prolonga a otro nivel; la fórmula no reinicia el estado en el espacio anterior.

La condición de memoria de esta proposición tiene una acción concreta en la cronología del TPK. Para una fase \(r\in\mathbb Z/9\mathbb Z\), con representante \(\bar r\in\{0,\ldots,8\}\), y una profundidad de memoria \(w\in\mathbb Z_{\geq0}\), cada emisión actualiza

\[
\sigma_9(w,r)=
\left(w+\left\lfloor\frac{\bar r+1}{9}\right\rfloor,r+[1]_9\right).
\]
Después de \(n\) emisiones, el número de cruces del representante \(8\) al representante \(0\) es \(\lfloor(\bar r+n)/9\rfloor\). Por tanto,

\[
\sigma_9^n(w,r)=
\left(w+\left\lfloor\frac{\bar r+n}{9}\right\rfloor,r+[n]_9\right),
\qquad
\sigma_9^9(w,r)=(w+1,r).
\]
Tomar \(q_k=w_k\) recupera el cociclo anterior y el incremento uno de cada vuelta. Incluso una emisión neutra en las hojas de APP registra una extensión en esta cronología. El cálculo no impone por sí mismo el operador espectral \(C\): determina la memoria sobre la que actúa la realización, una vez establecido su entrelazamiento.

\subsection{Calendario nonádico y contracción de una vuelta}
\label{rz:calendario}
El calendario finito es una representación de este transporte. Sobre \(\mathbb C^9\otimes\mathcal H_{\mathscr W}\), sea

\[
S_C=\sum_{r=1}^{8}|r+1\rangle\langle r|\otimes I
+|1\rangle\langle9|\otimes C.
\]
La unitariedad de \(C\) hace unitario a \(S_C\), y la vuelta completa satisface

\[
S_C^9=I_{\mathbb C^9}\otimes C.
\]
Cada vector atraviesa exactamente una vez la unión orientada que contiene \(C\); por eso no aparece \(C^9\).

Sea \(u=9^{-1/2}(1,\ldots,1)\) y \(\iota v=u\otimes v\). La compresión uniforme de un paso es

\[
A_C=\iota^*S_C\iota=\frac{8I+C}{9}.
\]
Un cálculo directo da

\[
I-A_C^*A_C
=\frac8{81}(I-C)^*(I-C).
\]
La diferencia respecto de un transporte unitario comprimido mide la componente que abandona el estado uniforme. No es una pérdida de unitariedad del calendario completo.

La contracción \(q=5/9\) de la hoja complementaria y el calendario son operaciones distintas. Si \(q\) corresponde a una vuelta nonádica, el operador por vuelta es

\[
M_{\mathrm{vuelta}}=qS_C^9
=q(I\otimes C).
\]
Puede repartirse uniformemente entre los nueve pasos mediante

\[
M_{\mathrm{paso}}=q^{1/9}S_C,\qquad
M_{\mathrm{paso}}^9=M_{\mathrm{vuelta}}.
\]
En cambio, tomar \(qS_C\) como operador de cada paso produce \(q^9(I\otimes C)\) al completar la vuelta. Las dos normalizaciones describen contracciones diferentes y no deben identificarse. Sobre una ordenada \(\gamma\), \(r\) vueltas del primer régimen producen

\[
q^r\tau_\gamma^r.
\]
El módulo registra la extinción del término terminal; la fase registra la acción sobre la memoria. Conocer sólo el módulo no determina las ordenadas espectrales.

\subsection{Matrices de momentos y aproximación espectral}
\label{rz:galerkin}
La realización anterior permite construir aproximaciones que parten de funciones de prueba, no de valores espectrales asignados. Elíjanse \(f_1,\ldots,f_d\in\mathcal D\) cuyas clases sean linealmente independientes en \(\mathcal H_{\mathscr W}\). Con \(e_j=[f_j]\), defínanse

\[
G_{ij}=\langle e_j,e_i\rangle,\qquad
K_{ij}=\langle He_j,e_i\rangle,\qquad
Q_{ij}=\langle He_j,He_i\rangle.
\]
La convención de índices hace que, para \(x=\sum_jv_je_j\), se tenga \(\|x\|^2=v^*Gv\). \(G\) es hermítica definida positiva y \(K,Q\) son hermíticas. Si las clases iniciales fueran dependientes, se debe retirar el núcleo de \(G\) antes de invertirla.

\HMTresultado{Proposición}{Aproximación de Galerkin y residuo}{rz-num-03-10}
El problema generalizado

\[
Kv=\gamma Gv
\]
equivale al problema hermítico ordinario de \(\widehat H=G^{-1/2}KG^{-1/2}\). Para cualquier \(\gamma\in\mathbb R\) y \(v\neq0\), el residuo espectral del vector correspondiente satisface

\[
\frac{\|(H-\gamma)x\|^2}{\|x\|^2}
=\frac{v^*(Q-2\gamma K+\gamma^2G)v}{v^*Gv}.
\]
En particular,

\[
\operatorname{dist}(\gamma,\sigma(H))
\leq\frac{\|(H-\gamma)x\|}{\|x\|}.
\]
\textbf{Demostración.} El cambio \(w=G^{1/2}v\) transforma la ecuación generalizada en \(\widehat Hw=\gamma w\). La expansión de la norma del residuo da la identidad cuadrática porque \(H\) es simétrico en los vectores utilizados. La última desigualdad se deduce de la medida espectral de \(H\):

\[
\|(H-\gamma)x\|^2
=\int |\lambda-\gamma|^2\,d\langle E_H(\lambda)x,x\rangle.
\]
El integrando está acotado inferiormente por la distancia al espectro al cuadrado. \hfill\(\square\)

La matriz de fase

\[
\widehat C=(\widehat H+iI/2)(\widehat H-iI/2)^{-1}
\]
es unitaria para el producto euclídeo. Si se utiliza en cambio \(G^{-1}K\) en la base original, la unitariedad correspondiente se expresa respecto de \(G\). Son dos bases del mismo problema, no dos normas intercambiables sin transformación.

Un residuo pequeño aproxima al espectro del operador ya construido. Por sí solo no cuenta todos sus autovalores, no demuestra la positividad global de \(\mathscr W\) y no certifica los decimales afectados por los errores de cuadratura. Para una cota validada, las entradas de \(G,K,Q\) deben acompañarse de cotas de error y de una cota inferior positiva de la forma \(v^*Gv\).

\subsection{Resolvente compacto y controles de exhaustividad}
\label{rz:resolvente}
En la realización atómica, las ordenadas forman un conjunto discreto y escapan a infinito. Por tanto,

\[
R=(H-iI)^{-1}
\]
es compacto: en la base atómica actúa por \((\gamma-i)^{-1}\), que tiende a cero. También se recupera directamente del grupo de traslaciones mediante

\[
R=i\int_0^\infty e^{-t}U_{-t}\,dt.
\]
La integral sobre un intervalo finito se define en sentido fuerte. Su cola desde \(T\) tiene norma operatoria a lo sumo \(e^{-T}\), por lo que las integrales truncadas convergen en norma. El cálculo funcional reduce la identidad a \(i/(1+i\gamma)=(\gamma-i)^{-1}\).

\HMTresultado{Proposición}{Convergencia en norma de las compresiones}{rz-num-03-11}
Si \(P_d\) son proyectores ortogonales de rango finito con \(P_d\to I\) fuertemente, entonces

\[
\|R-P_dRP_d\|\longrightarrow0.
\]
\textbf{Demostración.} La convergencia fuerte de \(P_d\), junto con la cota \(\|P_d\|\leq1\), es uniforme sobre todo compacto. Se aplica al cierre compacto de la imagen de la bola unidad por \(R\), lo que da \(\|(I-P_d)R\|\to0\). Aplicando el mismo argumento a \(R^*\) resulta \(\|R(I-P_d)\|\to0\). La descomposición

\[
R-P_dRP_d=(I-P_d)R+P_dR(I-P_d)
\]
concluye la prueba. \hfill\(\square\)

Esta convergencia permite controles por curvas de resolvente. Si una curva cerrada \(\mathcal C\) separa un grupo de autovalores no nulos y la norma del error de compresión es menor que el inverso de una cota del resolvente aproximado sobre \(\mathcal C\), la serie de Neumann conserva la invertibilidad a lo largo de la curva. Los proyectores de Riesz de los dos operadores tienen entonces el mismo rango. Es un control de exhaustividad en una región espectral delimitada, adicional al residuo de un vector.

La multiplicidad que cuenta ese rango es la del operador representado. En el espacio escalar de la subsección~\ref{rz:identificacion} no debe sustituirse por el peso \(m_\gamma\) sin introducir expresamente el observable ponderado.

\subsection{De las identidades internas al ensamblaje analítico}
\label{rz:alcance}
Se han reunido aquí las pruebas de la normalización de calor y Mellin, la construcción explícita de la forma completada, su descomposición por columnas, el transporte cilíndrico con memoria, los dos círculos, el calendario y los controles espectrales. La construcción aritmética de irreducibles precede a sus apariciones en los pesos primo–potencia. La memoria procede de las prolongaciones compatibles; no se sustituye por una cifra de fase.

La subsección~\ref{rz:columnas} reúne además un codificador interno finito construido con los transportes de historias, junto con la prueba de sus momentos \(I\) y \(T_N\) en la realización isométrica. La aplicación del teorema~\ref{rz-num-03-5} al codificador nativo de la forma completada conserva una identificación más específica: la acción previa \(R_R\) sobre funciones de prueba y la igualdad de aquellos momentos con los de \(J_{+,R}\) y \(J_{-,R}\). El propietario de la forma completada formula esas identidades y desarrolla sus consecuencias. Esta redacción conserva esas consecuencias con sus antecedentes explícitos y la composición interna efectivamente recuperada; no sustituye el enlace restante por una formulación nominal. La distinción se refiere a la demostración reunida en este artículo, no a una declaración de ausencia del resultado en el conjunto del corpus.

El enlace completo entre todas las historias APP–TRIT–TPK y la normalización por palabras del primer capítulo también conserva su obligación de integración. Las cinco construcciones consustanciales del continuo siguen perteneciendo al núcleo común; este capítulo expone su lectura espectral y no las redefine como estructuras independientes seleccionadas por problemas convencionales.

La identificación de la forma global con la fórmula explícita y el criterio de Weil fija exactamente el alcance de una eventual conclusión terminal. Las comprobaciones finitas del productor nativo y las aproximaciones espectrales son resultados verificables distintos. Por ello, esta edición de trabajo no presenta el capítulo como una demostración autónoma ya cerrada de la hipótesis de Riemann.

\subsection{Recuperación bidireccional y ensamblaje de los lectores}
\label{rz:gamma}
La información gamma recupera la prueba que la produjo. Esta propiedad proporciona una acción explícita entre las columnas completas, antes de conocer el signo de su diferencia. El desarrollo siguiente amplía el estado de integración descrito en las subsecciones~\ref{rz:columnas} y~\ref{rz:alcance}: el conector acotado global queda recuperado; su contractividad se estudia sobre la misma acción y el mismo rango, no sobre otro operador.

La construcción se sitúa después de los pesos primo–potencia, la traza completada y los transportes ya definidos. No cambia la selección de irreducibles ni recibe ceros como entradas. Los propietarios de la cola gamma, del dominio basal y del pliegue previo se conservan completos en el paquete de procedencia de esta integración.

\HMTresultado{Proposición}{Inversa izquierda de la cola y conector global}{rz-num-03-12}
Sea \(H_R=L^2(I_R)\), identificado con su extensión por cero a la recta, y sea \(M_R\) la restricción ortogonal a \(I_R\). Defínanse

\[
a_R=2\int_R^\infty\mu(a)\,da,\qquad
(B_Rf)(a)=\sqrt{\mu(a)}(I-T_a)f\quad(a\geq R).
\]
Entonces \(B_R\) es acotado y \(B_R^*B_R=a_RI\). Si \(\pi_{\Gamma,\geq R}\) extrae la cola gamma de la columna positiva, la aplicación

\[
L_R=a_R^{-1}B_R^*\pi_{\Gamma,\geq R}
\]
satisface \(L_RJ_{+,R}f=f\) para \(f\in\mathcal D_R\), y \(L_RL_R^*=a_R^{-1}I\) sobre \(H_R\). La primera identidad se extiende al dominio cerrado conjunto de los lectores; no afirma que la columna gamma completa esté definida en todo \(L^2(I_R)\). En particular,

\[
\mathcal C_R=J_{-,R}L_R,\qquad
\mathcal C_RJ_{+,R}=J_{-,R}
\]
es un conector acotado para todo \(R>0\).

\textbf{Demostración.} Para \(a\geq R\), los soportes de \(f,g\in H_R\) y de sus trasladadas están separados, salvo extremos de medida cero. Por tanto

\[
\langle(I-T_a)f,(I-T_a)g\rangle=2\langle f,g\rangle.
\]
Integrar da el primer Gram. El adjunto tiene la acción explícita

\[
B_R^*F=M_R\int_R^\infty\sqrt{\mu(a)}(I-T_{-a})F(a)\,da.
\]
La integral converge en norma por Cauchy–Schwarz; la cola de \(\mu\) es integrable y \(T_a^*=T_{-a}\). Aplicar el adjunto a la cola de \(J_{+,R}f\) devuelve \(a_Rf\). La extracción es una coisometría sobre su sumando, de donde \(L_RL_R^*=a_R^{-2}B_R^*B_R=a_R^{-1}I\). La columna negativa es acotada sobre \(H_R\): tiene un número finito de filas vectoriales y un funcional polar continuo en ese intervalo. La composición concluye. \hfill\(\square\)

La inversión usa las traslaciones conjugadas antes de expresar la función recuperada. El propietario proporciona también una versión de banda finita: con \(A_R=[R+1,R+2]\) y \(\nu_R=\int_{A_R}\mu(a)\,da\),

\[
\mathfrak R_RF=\nu_R^{-1}\int_{A_R}
\sqrt{\mu(a)}M_RF(a)\,da,\qquad
\mathfrak R_RD_\Gamma f=f.
\]
La identidad resulta de \(M_RT_af=0\) sobre esa banda. Las dos inversas son anteriores a una desigualdad de Weil. La primera permite calcular exactamente la norma ambiental de su conector.

\HMTresultado{Proposición}{Norma del conector y factorización basal}{rz-num-03-13}
Escribamos

\[
s_R^2=2\sum_{\ell_{p,k}\leq R}w_{p,k}-c_\Gamma,\qquad
\beta_R=2\sinh(R/2)-R.
\]
Entonces

\[
\|\mathcal C_R\|^2=\frac{s_R^2+\beta_R}{a_R}.
\]
Por consiguiente, si \(0<h<\log2\) y

\[
2M_\Gamma(h)\geq-c_\Gamma+2\sinh(h/2)-h,\qquad
M_\Gamma(h)=\int_h^\infty\mu(a)\,da,
\]
existe una dilatación que realiza los dos momentos del teorema~\ref{rz-num-03-5} sobre todas las funciones de \(\mathcal D_h\).

\textbf{Demostración.} La columna negativa tiene Gram \(J_-^*J_-=s_R^2I+b_-^*b_-\). El funcional

\[
b_-(f)=\sqrt2\int_{-R/2}^{R/2}f(x)\sinh(x/2)\,dx
\]
posee norma al cuadrado \(2\int_{-R/2}^{R/2}\sinh^2(x/2)\,dx=\beta_R\). El Gram tiene norma \(s_R^2+\beta_R\). La identidad anterior \(L_RL_R^*=a_R^{-1}I\) da \(\mathcal C_R\mathcal C_R^*=a_R^{-1}J_-J_-^*\), y prueba la fórmula.

Para \(R=h<\log2\) no hay longitudes primo–potencia en la suma. La desigualdad declarada prueba \(\|\mathcal C_h\|\leq1\). Se pueden construir, en este orden,

\[
D_h^{\mathrm{def}}=(I-\mathcal C_h^*\mathcal C_h)^{1/2},
\qquad
\Xi_hf=J_{-,h}f\oplus D_h^{\mathrm{def}}J_{+,h}f.
\]
La suma de los Grams de \(\mathcal C_h\) y de su defecto es la identidad. Como \(\mathcal C_hJ_{+,h}=J_{-,h}\), el primer momento de \(\Xi_h\) es \(J_{+,h}^*J_{+,h}\); la proyección al primer sumando tiene momento \(J_{-,h}^*J_{-,h}\). El teorema~\ref{rz-num-03-5} da la positividad. La raíz del defecto se toma después de demostrar contractividad, no como sustituto de esa prueba.

Si se amplía la ventana a \(R\geq h\), cada canal primo añadido cumple \(\|(I-T_\ell)f\|^2=2\|f\|^2\) sobre \(H_h\). La aplicación de esa imagen a \(\sqrt2f\) es una isometría, extendida por cero en el complemento de su imagen cerrada. La suma directa de esas filas con el conector gamma–polar mantiene la contractividad en el mismo dominio basal. \hfill\(\square\)

El extremo basal se determina por

\[
2M_\Gamma(h_*)=-c_\Gamma+2\sinh(h_*/2)-h_*.
\]
El cambio \(t=e^{-a/2}\) da

\[
M_\Gamma(h)=\operatorname{arctanh}(e^{-h/2})
+\arctan(e^{-h/2}).
\]
El primer miembro de la ecuación decrece estrictamente desde infinito hasta cero; el segundo es positivo y crece estrictamente para \(h>0\). Existe una única raíz positiva. El propietario la evalúa como \(h_*=0.08562601762973687068\ldots<\log2\). Puede trabajarse con cualquier \(h<h_*\) certificado por desigualdades, sin necesitar los decimales de ese extremo.

La forma completa también conserva la traslación de las pruebas. En los puertos polares ésta actúa mediante

\[
\binom{b_+(T_tf)}{b_-(T_tf)}=
\begin{pmatrix}\cosh(t/2)&\sinh(t/2)\\
\sinh(t/2)&\cosh(t/2)\end{pmatrix}
\binom{b_+(f)}{b_-(f)}.
\]
La matriz conserva la diferencia de cuadrados y las demás filas se trasladan unitariamente. El dominio basal puede por ello situarse en cualquier intervalo trasladado de longitud \(h\).

\HMTresultado{Proposición}{Unicidad sobre la imagen efectiva}{rz-num-03-14}
Sea \(\mathcal M_{+,R}=\overline{J_{+,R}\mathcal D_R}\). La restricción del conector construido a este subespacio es única entre las aplicaciones acotadas con esa identidad de salida. Además,

\[
\|\mathcal C_R|_{\mathcal M_{+,R}}\|\leq1
\quad\Longleftrightarrow\quad
\mathscr W_R(f,f)\geq0\quad(f\in\mathcal D_R).
\]
\textbf{Demostración.} Dos operadores acotados que coinciden en el rango denso coinciden en su cierre. La equivalencia resulta de \(\mathcal C_RJ_{+,R}f=J_{-,R}f\) y de la identidad de normas; la desigualdad se prolonga al cierre por continuidad. \hfill\(\square\)

La fórmula global está, por tanto, construida:

\[
\mathscr W_R(f,g)=
\langle J_{+,R}f,(I-\mathcal C_R^*\mathcal C_R)J_{+,R}g\rangle.
\]
Cambiar la representación ambiental puede facilitar el cálculo de la norma, pero no permite cambiar su acción sobre las pruebas. La no contractividad en todo el ambiente no refuta la positividad sobre la imagen completa, donde permanecen las correlaciones de las filas. La proyección a esa imagen tampoco prueba automáticamente la cota.

\HMTresultado{Proposición}{Pliegue previo y todos los productos cruzados}{rz-num-03-15}
Fijada una fibra basal, sean \(m=9^q\) y \(\ell_j\) los desplazamientos expresados por sus rutas. Sobre \(m\) copias de \(L^2(\mathbb R)\), el pliegue y su representación uniforme son

\[
\mathcal F_qx=\sum_{j=1}^mT_{\ell_j}x_j,\qquad
\mathsf S_qx=u_q\otimes\mathcal F_qx=3^qP_q\mathsf T_qx,
\qquad u_q=m^{-1/2}(1,\ldots,1).
\]
Para todo lector lineal \(J\) definido sobre esas pruebas,

\[
\|J\mathcal F_qx\|^2
=\sum_{i,j}\langle JT_{\ell_i}x_i,JT_{\ell_j}x_j\rangle.
\]
Además, \(\mathcal F_q\) es sobreyectiva, con inversa derecha

\[
(\mathcal E_qf)_j=m^{-1}T_{-\ell_j}f,\qquad
\mathcal F_q\mathcal E_q=I,\qquad
\|\mathcal E_qf\|^2=m^{-1}\|f\|^2.
\]
\textbf{Demostración.} Promediar las \(m\) rutas y multiplicar por \(\sqrt m=3^q\) da la expresión de \(\mathsf S_q\). Expandir el producto interior conserva todos los cruces. Sustituir \(\mathcal E_qf\) produce \(m\) copias de \(f/m\); la unitariedad de las traslaciones da la norma. \hfill\(\square\)

Si los sucesores cumplen \(\ell_{jr}=\ell_j+\delta_{jr}\), la operación \((\mathcal A_qx)_j=\sum_{r=0}^8T_{\delta_{jr}}x_{jr}\) satisface exactamente

\[
\mathcal F_{q+1}=\mathcal F_q\mathcal A_q,\qquad
\mathsf S_{q+1}=\mathcal I_q\mathsf S_q\mathcal A_q,\qquad
\mathcal I_qy=y\otimes u_9.
\]
Las columnas gamma, primo y polar se aplican así a la misma suma antes de sus normas. Una mezcla unitaria posterior a una suma ortogonal preserva la energía diagonal; no fabrica cruces omitidos.

La sobreyectividad determina el alcance del refinamiento. En la fibra completa, la positividad de \(\mathscr W(\mathcal F_qx,\mathcal F_qx)\) para todos los estados equivale a la positividad de \(\mathscr W(f,f)\) para todas las pruebas. La implicación inversa usa \(x=\mathcal E_qf\), cuyas componentes siguen siendo suaves de soporte compacto. El pliegue preserva el problema y sus cruces: no restringe las pruebas por el solo hecho de aumentar las rutas. Una restricción a estados alcanzables más específicos exige su construcción y su prueba.

\subsubsection*{Resultado reunido y extensión global}
La inversa gamma, el enlace completo de salida, el dominio basal, su traslación y el pliegue anterior a los lectores tienen aquí acción, dominio y demostración. Son resultados recuperados del corpus y formalizaciones reunidas, no descubrimientos atribuidos retrospectivamente.

Para el alcance global solicitado resta demostrar la cota de esa misma salida sobre toda la imagen efectiva, incluidos los cruces de las sumas de pruebas locales, de manera compatible al ampliar las ventanas. La identidad de salida de un conector y la contractividad de otro no se combinan sin demostrar que representan la misma acción. La comparación con el codificador interno de historias conserva los dos momentos especificados en la subsección~\ref{rz:columnas}.

\subsection{Ventanas de relojes irreducibles y refinamiento con términos cruzados}
\label{rz:ventanas}
El producto nativo determina los irreducibles y sus potencias antes de la lectura logarítmica. Las longitudes \(\ell_{p,k}=k\log p\) y los pesos \(w_{p,k}=(\log p)p^{-k/2}\) que aparecen a continuación son los ya expresados en las secciones~\ref{lectura:manuscrito-01} y~\ref{lectura:manuscrito-02}. No se seleccionan desde el signo que se quiere comprobar. Sobre esta salida aritmética puede calcularse conjuntamente una colección de pruebas, manteniendo los canales de la forma completada y todos sus productos cruzados.

Esta sección reúne la densidad y el criterio de Schur del propietario «Lector superviviente de dos rutas APP», y extiende el cálculo a intervalos trasladados de igual anchura. La elección de nueve ventanas que se certifica abajo es una colección analítica explícita; su cardinal no identifica por sí solo esas ventanas con los nueve sucesores canónicos de un estado TPK.

\HMTresultado{Proposición}{Núcleo completo de intervalos trasladados}{rz-num-03-16}
Para \(\varepsilon>0\), sean

\[
b_{\varepsilon,t}(x)=\varepsilon^{-1/2}
\mathbf1_{[t-\varepsilon/2,t+\varepsilon/2]}(x),\qquad
\tau_\varepsilon(x)=\left(1-\frac{|x|}{\varepsilon}\right)_+,
\]
\[
a_\varepsilon=\frac{4\sinh(\varepsilon/4)}{\sqrt\varepsilon},
\qquad \lambda_n=2n+\frac12,\qquad
F(x)=\sum_{n\geq0}\frac{e^{-\lambda_nx}}{\lambda_n^2}
\quad(x\geq0).
\]
La forma completa sobre estos intervalos depende de la diferencia de centros:

\[
\mathscr W(b_{\varepsilon,s},b_{\varepsilon,t})
=K_\varepsilon(s-t),
\]
donde

\[
\begin{aligned}
K_\varepsilon(d)={}&2a_\varepsilon^2\cosh(d/2)
+c_\Gamma\tau_\varepsilon(d)\\
&+\frac{2F(|d|)-F(|d-\varepsilon|)-F(|d+\varepsilon|)}{\varepsilon}\\
&-\sum_{p,k}w_{p,k}\bigl[
\tau_\varepsilon(d-\ell_{p,k})+
\tau_\varepsilon(d+\ell_{p,k})\bigr].
\end{aligned}
\]
La suma primo–potencia es finita; sólo contribuyen longitudes \(\ell_{p,k}<|d|+\varepsilon\).

\textbf{Demostración.} El solapamiento de dos intervalos normalizados es la correlación triangular \(\tau_\varepsilon\). Al insertar las traslaciones de cada reloj se obtienen los dos términos primos. Las integrales polares de \(b_{\varepsilon,t}\) son \(a_\varepsilon e^{t/2}\) y \(a_\varepsilon e^{-t/2}\). Su producto cruzado da \(a_\varepsilon^2(e^{(s-t)/2}+e^{(t-s)/2})\), el término polar de la fórmula. En particular, la dependencia de la diferencia de centros pertenece a la forma completa, aunque cada puerto polar por separado cambie bajo traslación.

La expansión \(\mu(a)=\sum_{n\geq0}e^{-\lambda_na}\) y una integración en los tramos lineales de \(\tau_\varepsilon\) dan

\[
\int_0^\infty e^{-\lambda a}
\bigl[2\tau_\varepsilon(d)-\tau_\varepsilon(d-a)
-\tau_\varepsilon(d+a)\bigr]\,da
=\frac{2e^{-\lambda|d|}-e^{-\lambda|d-\varepsilon|}
-e^{-\lambda|d+\varepsilon|}}{\varepsilon\lambda^2}.
\]
Para justificar el intercambio de serie e integral, la función triangular cumple además

\[
|2\tau_\varepsilon(d)-\tau_\varepsilon(d-a)-\tau_\varepsilon(d+a)|
\leq\min(4,2a/\varepsilon).
\]
Su producto por \(\mu(a)\) es integrable tanto en el origen como en infinito. La suma de las integrales absolutas está así acotada, lo que justifica el intercambio. El sumando integrado queda también acotado en valor absoluto por \(4/(\varepsilon\lambda_n^2)\), cuya serie converge. La intervención conjunta de los dos lados cancela la singularidad de \(\mu\) en el origen. Se puede, por tanto, sumar la identidad integrada y obtener el término gamma declarado. \hfill\(\square\)

Estos indicadores son elementos del dominio cerrado de los lectores, no funciones suaves. La extensión desde \(\mathcal D_R\) se justifica sin modificar la forma. Para un intervalo normalizado,

\[
\|b_{\varepsilon,t}-T_ab_{\varepsilon,t}\|_2^2
=2\min(|a|/\varepsilon,1).
\]
La integral contra \(\mu(a)\) es finita. Sea \(b_{\varepsilon,t}^{(\delta)}=\rho_\delta*b_{\varepsilon,t}\) una regularización suave, de soporte contenido en una ventana fija. Se toma \(\rho\geq0\), suave y compactamente soportada, con \(\int\rho=1\), y \(\rho_\delta(x)=\delta^{-1}\rho(x/\delta)\). El símbolo gamma es

\[
\sigma_\Gamma(\xi)=\int_0^\infty
\mu(a)|1-e^{-ia\xi}|^2\,da.
\]
En Fourier, el integrando que expresa la norma gamma de la diferencia contiene \(\sigma_\Gamma|\widehat b|^2 |\widehat\rho(\delta\xi)-1|^2\). El primer producto es integrable y el último factor está acotado por \(4\) y tiende a cero. La convergencia dominada prueba la convergencia gamma. Los canales primos son finitos y acotados en esa ventana; los polares son funcionales continuos sobre su soporte. Se obtiene convergencia en la norma conjunta de los lectores y, en consecuencia, de cada Gram finito.

\HMTresultado{Corolario}{Dos orientaciones de una misma ventana prima}{rz-num-03-17}
Póngase \(D_\varepsilon=K_\varepsilon(0)\). Para dos centros separados por \(L\), el Gram completo es

\[
G_2=\begin{pmatrix}D_\varepsilon&K_\varepsilon(L)\\
K_\varepsilon(L)&D_\varepsilon\end{pmatrix}.
\]
Por tanto

\[
\mathscr W(af+bg,af+bg)=
\frac{D_\varepsilon+K_\varepsilon(L)}2|a+b|^2
+\frac{D_\varepsilon-K_\varepsilon(L)}2|a-b|^2.
\]
La positividad sobre su espacio generado exige simultáneamente \(D_\varepsilon+K_\varepsilon(L)\geq0\) y \(D_\varepsilon-K_\varepsilon(L)\geq0\).

Si \(L\geq\varepsilon\), defínanse

\[
I_\varepsilon(L)=
\frac{F(L-\varepsilon)-2F(L)+F(L+\varepsilon)}{\varepsilon}>0,
\]
\[
P_\varepsilon(L)=
\sum_{|\ell_{p,k}-L|<\varepsilon}w_{p,k}
\left(1-\frac{|\ell_{p,k}-L|}{\varepsilon}\right).
\]
La convexidad estricta de cada exponencial prueba la positividad de \(I_\varepsilon\). La condición de los dos signos se escribe exactamente

\[
2a_\varepsilon^2\cosh(L/2)-I_\varepsilon(L)-D_\varepsilon
\leq P_\varepsilon(L)
\leq
2a_\varepsilon^2\cosh(L/2)-I_\varepsilon(L)+D_\varepsilon.
\]
Así, el cruce completo controla una ventana de longitudes de relojes irreducibles con sus pesos, y no sólo un valor diagonal. La identidad resulta de la proposición~\ref{rz-num-03-16} y de diagonalizar la matriz de dos entradas; no presupone su signo. \hfill\(\square\)

\HMTresultado{Proposición}{Control conjunto de nueve pruebas}{rz-num-03-18}
Sean \(\varepsilon=1/100\), \(t_j=j\log2\), \(0\leq j\leq8\), y \(G_{ij}=K_\varepsilon(t_i-t_j)\). El cálculo por intervalos conservado con este capítulo certifica

\[
\lambda_{\min}(G)>0.8479.
\]
En consecuencia,

\[
\mathscr W\left(\sum_{j=0}^8z_jb_{\varepsilon,t_j},
\sum_{j=0}^8z_jb_{\varepsilon,t_j}\right)
>0.8479\sum_{j=0}^8|z_j|^2
\quad(z\neq0).
\]
\textbf{Demostración y reproducción.} Para \(F_N(x)=\sum_{n=0}^{N-1}e^{-\lambda_nx}/\lambda_n^2\), la función sumada decrece y

\[
0\leq F(x)-F_N(x)
\leq e^{-\lambda_Nx}
\left(\lambda_N^{-2}+\frac1{2\lambda_N}\right).
\]
En efecto, se separa el primer término de la cola y se acota el resto por la integral de \((2u+1/2)^{-2}\), extrayendo la exponencial mayor. Se usan \(N=1024\) y operaciones intervalares con redondeo hacia fuera. La condición \(8\log2+\varepsilon<\log259\) permite censar todos los pares primo–potencia hasta \(259\): resultan \(71\) filas. Los irreducibles y productos provienen del productor nativo ya conservado; el reconocimiento por enteros ordinarios es posterior. La constante \(c_\Gamma\) se evalúa desde los intervalos certificados de las coordenadas anteriores, no desde un valor objetivo de esta forma.

La fórmula de la proposición~\ref{rz-num-03-16} da intervalos para las nueve entradas de la primera fila. Para cada fila \(i\), se verifica con dichos intervalos

\[
G_{ii}-\sum_{j\neq i}|G_{ij}|>0.8479.
\]
La matriz es hermítica. La desigualdad \(2|z_i||z_j|\leq|z_i|^2+|z_j|^2\), aplicada a cada cruce, da directamente la cota cuadrática anunciada. La verificación conserva los intervalos de las nueve filas y el censo completo. Los márgenes estrictos y la convergencia de Grams demostrada antes permiten también regularizar simultáneamente las nueve pruebas, obteniendo positividad sobre su espacio suave generado para una regularización suficientemente fina. \hfill\(\square\)

Este control cubre todas las combinaciones complejas de nueve pruebas, no únicamente la suma de ellas. Su dimensión permanece finita. El registro computacional declara ese alcance y no usa su resultado como certificado de una desigualdad universal.

\HMTresultado{Proposición}{Refinamiento y ocho componentes de detalle}{rz-num-03-19}
La subdivisión geométrica de un intervalo en nueve partes cumple

\[
b_{\varepsilon,t}=\frac13\sum_{r=0}^8
b_{\varepsilon/9,t+(r-4)\varepsilon/9}.
\]
Si \(V=(1,\ldots,1)^{\mathsf T}/3\), entonces

\[
G_{\mathrm{padre}}=V^*G_{\mathrm{hijos}}V.
\]
Para varios predecesores se toma la inclusión isométrica por bloques con esa misma columna. En una base formada por la componente uniforme y las ocho componentes de detalle por predecesor, escribamos

\[
G_{\mathrm{hijos}}=
\begin{pmatrix}A&B^*\\B&D\end{pmatrix},\qquad
A=G_{\mathrm{padres}}.
\]
Si \(A>0\), la condición exacta es

\[
G_{\mathrm{hijos}}\succeq0
\quad\Longleftrightarrow\quad D-BA^{-1}B^*\succeq0.
\]
\textbf{Demostración.} Los nueve intervalos sucesores particionan al predecesor salvo extremos de medida cero. La normalización de sus indicadores prueba la identidad lineal, y la sesquilinealidad da la identidad de Grams con todos los cruces. Finalmente,

\[
\begin{aligned}
\begin{pmatrix}x\\y\end{pmatrix}^*
\begin{pmatrix}A&B^*\\B&D\end{pmatrix}
\begin{pmatrix}x\\y\end{pmatrix}
={}&(x+A^{-1}B^*y)^*A(x+A^{-1}B^*y)\\
&+y^*(D-BA^{-1}B^*)y.
\end{aligned}
\]
Tomar \(x=-A^{-1}B^*y\) prueba la necesidad y la misma identidad prueba la suficiencia. \hfill\(\square\)

La componente uniforme reproduce el Gram anterior; los detalles contienen las nuevas interferencias. Conservar ambas partes es la exigencia exacta para prolongar una prueba positiva bajo este refinamiento. El retorno de nueve fases no convierte por sí solo el Gram en una matriz circulante: en la colección anterior \(K_\varepsilon(\log2)\neq K_\varepsilon(8\log2)\). La fase y la memoria de distancias siguen siendo datos distintos.

La matriz de detalle y su complemento de Schur proporcionan una condición comprobable para cada ampliación. No se define el lector prospectivo mediante una factorización retrospectiva de la forma; se comprueba la forma expresada por los lectores previamente construidos. Para completar el argumento global se requiere una estimación sobre todos esos detalles, uniforme en las ampliaciones cofinales y en el dominio completo de pruebas.

\subsection{Ensamblaje coercivo y estabilidad de los detalles nonádicos}
\label{rz:ensamblaje}
La propiedad basal permite controlar espacios completos de funciones. Para ensamblar varios intervalos hay que comparar ese margen local con cada interacción entre rutas. El núcleo completo de la sección anterior proporciona una estimación de esas interacciones antes de elegir una función. Se obtiene así un resultado de dimensión infinita sobre un dominio desconectado, no sólo un Gram de indicadores.

Las posiciones se toman después de la representación de los relojes irreducibles. El caso explícito utiliza las traslaciones sucesivas del reloj nativo 2 y una anchura posicional declarada. Los operadores de diferencia conservan las dos orientaciones. El refinamiento posterior subdivide cada intervalo en nueve partes y mantiene la misma función antes de sus dos lecturas de energía. Esta colección analítica se especifica como tal: su cardinal no la identifica por sí solo con una selección canónica de estados TPK.

\HMTresultado{Proposición}{Margen diagonal y cota de los cruces}{rz-num-03-20}
Sean \(t_1<\cdots<t_m\) centros distintos e \(I_j=(t_j-\varepsilon/2,t_j+\varepsilon/2)\) intervalos disjuntos, con \(0<\varepsilon<\log2\). Defínase

\[
\delta_\varepsilon=2M_\Gamma(\varepsilon)+c_\Gamma
-2\sinh(\varepsilon/2)+\varepsilon.
\]
Para \(f_j\in C_c^\infty(I_j)\),

\[
\mathscr W(f_j,f_j)\geq\delta_\varepsilon\|f_j\|_2^2.
\]
Si \(d_{ij}=|t_i-t_j|>\varepsilon\), póngase

\[
\rho_{ij}(\varepsilon)=
\sum_{|\ell_{p,k}-d_{ij}|<\varepsilon}w_{p,k}
+\varepsilon\mu(d_{ij}-\varepsilon)
+2\varepsilon\cosh((d_{ij}+\varepsilon)/2).
\]
Entonces

\[
|\mathscr W(f_i,f_j)|
\leq\rho_{ij}(\varepsilon)\|f_i\|_2\|f_j\|_2
\qquad(i\neq j).
\]
\textbf{Demostración.} En un intervalo de anchura \(\varepsilon\), la cola gamma aporta al menos \(2M_\Gamma(\varepsilon)\|f_j\|_2^2\). No hay correlación primo–potencia diagonal, porque su longitud mínima es \(\log2>\varepsilon\). En una ventana ambiental mayor, cada diferencia prima y su contratermino se cancelan sobre ese soporte. Tras centrar el intervalo, el puerto polar negativo tiene norma al cuadrado \(2\sinh(\varepsilon/2)-\varepsilon\). Descartar el puerto polar positivo y la parte gamma anterior a la cola da la cota. La invariancia por traslación de la forma completa lleva la misma cota a cualquier centro.

En el cruce, la contribución escalar es nula por disjunción de los soportes. Para cada longitud positiva, sólo una de las dos orientaciones puede conectar ambos intervalos. La unitariedad de la traslación acota la contribución prima por la suma de pesos indicada. El cruce gamma tiene núcleo \(-\mu(|x-y|)\) sobre \(I_i\times I_j\). Como \(\mu\) es decreciente,

\[
|\mathscr W_\Gamma(f_i,f_j)|
\leq\mu(d_{ij}-\varepsilon)\|f_i\|_1\|f_j\|_1
\leq\varepsilon\mu(d_{ij}-\varepsilon)\|f_i\|_2\|f_j\|_2.
\]
El núcleo polar completo es \(2\cosh((x-y)/2)\). La cota \(|x-y|\leq d_{ij}+\varepsilon\) y la misma desigualdad entre las normas \(L^1\) y \(L^2\) dan el término polar. Sumar las tres estimaciones prueba el resultado. \hfill\(\square\)

\HMTresultado{Teorema}{Coercividad del ensamblaje}{rz-num-03-21}
Si

\[
\eta=\delta_\varepsilon-
\max_i\sum_{j\neq i}\rho_{ij}(\varepsilon)>0,
\]
entonces, para \(\Omega=\bigcup_j I_j\),

\[
\boxed{\mathscr W(f,f)\geq\eta\|f\|_2^2
\quad\text{para toda }f\in C_c^\infty(\Omega).}
\]
La afirmación vale en el cierre de ese dominio para la norma conjunta de los lectores. En particular, para

\[
m=9,\qquad t_j=j\log2\ (0\leq j\leq8),\qquad
\varepsilon=1/1000,
\]
la evaluación certificada da \(\eta>1\).

\textbf{Demostración.} Escribiendo \(f=\sum_jf_j\) sobre los soportes disjuntos y \(x_j=\|f_j\|_2\), se obtiene

\[
\begin{aligned}
\mathscr W(f,f)&\geq
\delta_\varepsilon\sum_jx_j^2
-2\sum_{i<j}\rho_{ij}(\varepsilon)x_ix_j\\
&\geq\sum_i\left(\delta_\varepsilon-
\sum_{j\neq i}\rho_{ij}(\varepsilon)\right)x_i^2.
\end{aligned}
\]
Se ha usado \(2x_ix_j\leq x_i^2+x_j^2\) para cada cruce. La suma de las normas al cuadrado es \(\|f\|_2^2\). La continuidad de los lectores extiende la desigualdad al dominio cerrado.

En el caso explícito, las distancias son \(d\log2\), \(1\leq d\leq8\). Para \(\varepsilon=1/1000\),

\[
2^d(e^\varepsilon-1)<\frac{256}{999}<1,
\qquad 2^d(1-e^{-\varepsilon})<\frac{256}{1000}<1.
\]
La primera desigualdad usa la serie exponencial y \(e^\varepsilon<1/(1-\varepsilon)\); la segunda usa \(1-e^{-\varepsilon}<\varepsilon\). Por integridad de las potencias expresadas, el único reloj primo–potencia en la ventana de esa distancia es \((p,k)=(2,d)\). Luego

\[
\rho_d=(\log2)2^{-d/2}
+\frac1{1000}\mu(d\log2-1/1000)
+\frac2{1000}\cosh((d\log2+1/1000)/2).
\]
El productor nativo verifica además el censo completo hasta \(259\). El cálculo de \(\delta_\varepsilon\) utiliza la fórmula ya demostrada de \(M_\Gamma\) y la identidad

\[
\arctan t=\frac\pi4-
\arctan\frac{1-t}{1+t}\qquad(0<t<1).
\]
El segundo argumento es pequeño y su serie alternada lleva la cota del primer término omitido. La representación interna de \(\pi\), el intervalo racional de gamma y las operaciones intervalares posteriores producen las cotas

\[
\delta_\varepsilon>4.4921,\qquad
\max_i\sum_{j\neq i}\rho_{ij}<2.5385,\qquad
\eta>1.9536>1.
\]
El recibo conserva el margen de cada fila. No se evalúan funciones particulares para inferir la desigualdad universal sobre \(C_c^\infty(\Omega)\): ésta procede de la estimación operatoria anterior. \hfill\(\square\)

\HMTresultado{Corolario}{Todos los refinamientos del dominio ensamblado}{rz-num-03-22}
Subdivídase cada \(I_j\) sucesivamente en nueve partes. Los indicadores normalizados de los \(9^{n+1}\) intervalos de profundidad \(n\) forman una colección ortonormal en \(L^2\). Su Gram completo satisface

\[
G_n\succeq\eta I\qquad(n\geq0).
\]
La inclusión uniforme \(V_n\), con nueve componentes \(1/3\) por predecesor, cumple

\[
V_n^*G_{n+1}V_n=G_n.
\]
En la descomposición uniforme–detalle,

\[
G_{n+1}=\begin{pmatrix}A_n&B_n^*\\B_n&D_n\end{pmatrix},
\qquad A_n=G_n,
\]
se obtiene, a cualquier profundidad,

\[
\boxed{D_n-B_nA_n^{-1}B_n^*\succeq\eta I.}
\]
\textbf{Demostración.} El teorema se aplica a las regularizaciones de cada combinación de indicadores. La convergencia en norma de los lectores de la subsección~\ref{rz:ventanas} permite pasar al límite. Un indicador que alcanza un extremo de \(I_j\) se aproxima primero contrayendo ligeramente su soporte hacia el centro, y luego regularizando dentro de \(I_j\). La correlación triangular y su cota integrable de traslación prueban la misma convergencia; no se utiliza una extensión del dominio de coercividad. La ortonormalidad en \(L^2\) transforma la cota funcional en \(G_n\succeq\eta I\), y la subdivisión exacta da la identidad de compresión.

Para un vector de detalle \(y\), se toma \(x=-A_n^{-1}B_n^*y\). La identidad de Schur da

\[
\begin{aligned}
y^*(D_n-B_nA_n^{-1}B_n^*)y
&=\begin{pmatrix}x\\y\end{pmatrix}^*G_{n+1}
\begin{pmatrix}x\\y\end{pmatrix}\\
&\geq\eta(\|x\|^2+\|y\|^2)\geq\eta\|y\|^2.
\end{aligned}
\]
La misma \(\eta\) sirve a cualquier profundidad. \hfill\(\square\)

El corolario controla los ocho detalles nuevos por predecesor a profundidad arbitraria dentro del ensamblaje. La función no se sustituye por su media ni se descartan sus componentes de suma nula. Es un fortalecimiento del Gram de nueve indicadores: cubre espacios funcionales de dimensión infinita y da una cota independiente de la profundidad.

\subsubsection*{Colecciones finitas y extensión del dominio}
Para cualquier conjunto finito de centros distintos, existe una anchura positiva suficientemente pequeña que permite aplicar el teorema~\ref{rz-num-03-21}. En efecto, \(2M_\Gamma(\varepsilon)\to+\infty\) cuando \(\varepsilon\downarrow0\), mientras que el término polar negativo local tiende a cero. Las distancias entre centros están separadas de cero. Las sumas de pesos primos quedan acotadas al reducir cada ventana, porque hay un número finito de longitudes bajo cualquier cota fijada. Los términos cruzados gamma y polar tienden a cero. Por tanto, \(\delta_\varepsilon\) acaba superando la suma de interacciones de cada fila.

La anchura admisible depende de los centros. Las uniones de intervalos estrechos obtenidas de esta manera no forman, por la sola construcción anterior, una cobertura cofinal de todas las funciones de soporte compacto. La positividad global requiere controlar también el crecimiento del dominio y sus nuevas interacciones. Quedan diferenciadas la uniformidad demostrada en profundidad y la extensión a cualquier dominio de pruebas, no demostrada aquí.

\subsection{Coercividad de los detalles en una ventana arbitraria fijada}
\label{rz:detalles}
La estimación precedente conserva un margen uniforme al refinar el interior de nueve intervalos separados. Una segunda construcción permite controlar los detalles finos de cualquier ventana fijada, incluidas sus celdas contiguas. Ambas afirmaciones tienen dominios diferentes y complementarios: la primera estima la función completa en un soporte desconectado; la segunda estima las componentes de media nula por celda en una ventana completa.

Se mantienen los relojes, pesos y lectores construidos previamente. En particular, las longitudes \(\ell_{p,k}\), los pesos \(w_{p,k}\), el escalar \(c_\Gamma\) y la representación interna \(\pi\) conservan su procedencia. El argumento que sigue actúa sobre la realización funcional de esos objetos; no modifica su generador. Escríbase

\[
\beta_R=2\sinh(R/2)-R,\qquad
\mathfrak c_R=\frac23-c_\Gamma+
2\sum_{\ell_{p,k}\leq R}w_{p,k}+\beta_R.
\]
La suma es finita y procede del producto nativo bajo la cota \(e^R\). Divídase \(I_R=(-R/2,R/2)\) en \(M\) celdas de longitud a lo sumo \(h\), con \(0<h\leq1\).

\HMTresultado{Teorema}{Umbral explícito para los detalles de media nula}{rz-num-03-23}
Sea \(d\) una función del dominio conjunto de los lectores, soportada en \(I_R\), cuya integral en cada celda es cero. Entonces

\[
\boxed{\mathscr W_R(d,d)\geq
\bigl[\log(1+4/h)-\mathfrak c_R\bigr]\|d\|_2^2.}
\]
En consecuencia, la forma es estrictamente coerciva en ese espacio de detalles si

\[
0<h<\min\left(1,\frac4{\exp(\mathfrak c_R)-1}\right).
\]
El umbral depende de la ventana \(R\), pero no de la función ni de una enumeración finita de ejemplos.

\textbf{Demostración.} En cada celda \([a,b]\) se define \(F(x)=\int_a^x d(t)\,dt\). La condición de media nula implica \(F(a)=F(b)=0\). Al reunir las primitivas se obtiene \(F\in H_0^1(I_R)\), extendida por cero fuera de la ventana, con \(F'=d\). La desigualdad de Dirichlet en cada celda da

\[
\|F\|_2^2\leq\frac{h^2}{\pi^2}\|d\|_2^2.
\]
Aquí \(\pi\) es la misma coordenada interna reconocida en la realización funcional. No se selecciona una nueva constante.

La cola gamma admite la expansión positiva ya utilizada

\[
\mu(a)=\sum_{n=0}^{\infty}e^{-\lambda_n a},
\qquad \lambda_n=2n+\frac12.
\]
Para \(\lambda>0\), sea \(C_\lambda\) la convolución por \(e^{-\lambda|x|}\). El aporte de un sumando es

\[
E_\lambda(d)=\int_0^\infty e^{-\lambda a}
\|(I-T_a)d\|_2^2\,da
=\frac2\lambda\|d\|_2^2-
\langle d,C_\lambda d\rangle\geq0.
\]
Su positividad es la de una integral de normas al cuadrado. Denótese por \(\mathcal F_u F=(2\pi)^{-1/2}\widehat F\) la transformada de Fourier unitaria, distinguiéndola de la transformada no normalizada empleada antes. El multiplicador de \(C_\lambda\) es \(2\lambda/(\lambda^2+\xi^2)\). Como \(d=F'\), Plancherel implica

\[
\begin{aligned}
\langle d,C_\lambda d\rangle
&=\int_{\mathbb R}
\frac{2\lambda\xi^2}{\lambda^2+\xi^2}
|(\mathcal F_u F)(\xi)|^2\,d\xi\\
&\leq2\lambda\|F\|_2^2
\leq\frac{2\lambda h^2}{\pi^2}\|d\|_2^2.
\end{aligned}
\]
Por Tonelli, la forma gamma es la suma de los \(E_{\lambda_n}\). Se conservan los primeros \(N+1\) términos y se descarta únicamente una cola no negativa. Resulta

\[
\begin{aligned}
Q_\Gamma(d)&\geq S_N(h)\|d\|_2^2,\\
S_N(h)&=4\sum_{n=0}^N\frac1{4n+1}
-\frac{h^2(N+1)(2N+1)}{\pi^2}.
\end{aligned}
\]
Elíjase \(N=\lfloor1/h\rfloor\). Entonces \(\lambda_N h\leq2+h/2<\pi\); las cotas retenidas son no negativas. La comparación con la integral de una función decreciente da

\[
4\sum_{n=0}^N\frac1{4n+1}
\geq\int_0^{N+1}\frac4{4x+1}\,dx
=\log(4N+5)\geq\log(1+4/h).
\]
Además, \(\pi^2>9\) y \(h\leq1\) implican

\[
\frac{h^2(N+1)(2N+1)}{\pi^2}
\leq\frac{(1+h)(2+h)}{\pi^2}<\frac23.
\]
Por tanto, \(Q_\Gamma(d)\geq[\log(1+4/h)-2/3]\|d\|_2^2\). En la forma completa, el término escalar aporta \(c_\Gamma\|d\|_2^2\); la parte prima se acota inferiormente por \(-2\sum_{\ell_{p,k}\leq R}w_{p,k}\|d\|_2^2\) mediante la unitariedad de cada traslación; y la parte polar se acota por \(-\beta_R\|d\|_2^2\), como se demostró en la subsección~\ref{rz:gamma}. Su suma prueba la desigualdad anunciada. El umbral escrito equivale a \(\log(1+4/h)>\mathfrak c_R\). El paso al dominio cerrado conserva la estimación por continuidad de los lectores y de las integrales de celda. \hfill\(\square\)

\HMTresultado{Corolario}{Dimensión finita de las posibles direcciones negativas}{rz-num-03-24}
Fijados \(R\) y una partición que cumpla el umbral anterior, cualquier subespacio del dominio en el que \(\mathscr W_R\) sea estrictamente negativa definida tiene dimensión a lo sumo \(M\), el número de celdas.

\textbf{Demostración.} Considérese la aplicación lineal

\[
\mathcal P(d)=\left(\int_{I_1}d,\ldots,
\int_{I_M}d\right)\in\mathbb C^M.
\]
Su núcleo es el espacio de detalles del teorema~\ref{rz-num-03-23}, sobre el que la forma es positiva definida. La restricción de \(\mathcal P\) a un subespacio estrictamente negativo es inyectiva: un vector no nulo de su núcleo tendría a la vez valor positivo y negativo. La dimensión de ese subespacio no puede superar \(M\). \hfill\(\square\)

Esta reducción conserva una cuestión finita de acoplamiento en cada ventana. No equivale a probar que el índice negativo sea cero. En una descomposición media–detalle, el bloque de detalles es coercivo a suficiente profundidad, pero deben conservarse el bloque de medias y los términos cruzados. Si se utiliza el complemento de Schur, su signo requiere esa composición completa. En particular, una resonancia prima exacta entre dos celdas no desaparece por imponer media nula: puede actuar como un múltiplo negativo de la identidad sobre sus detalles correspondientes. El teorema mantiene tales términos mediante su cota de traslación.

El resultado aporta una segunda uniformidad efectiva: en cualquier ventana fijada, los detalles suficientemente finos tienen una cota común explícita. La subsección~\ref{rz:ensamblaje} prueba además el signo del ensamblaje completo, con sus cruces, para la unión de nueve intervalos especificada. La extensión de ese signo a toda ventana sigue siendo el paso global que esta edición no da por demostrado.
